% =============================================================
% Capitulo 7 - Diseno e implementacion de la placa de control
% =============================================================
\chapter{Dise\~no e implementaci\'on de la placa de control}
\label{cap:placa_control}

\section{Selecci\'on del microcontrolador y arquitectura general}
\label{sec:arquitectura_general}

La placa de control integra un microcontrolador RP2350, sobre el
m\'odulo comercial Raspberry Pi Pico, seleccionado seg\'un los criterios
descritos en la Secci\'on~\ref{sec:rp2350} \cite{informeEneFeb2026,kicad_hsrl}.
A nivel de arquitectura general, la placa recibe la se\~nal
acondicionada por el circuito detector de picos (Cap\'itulo~\ref{cap:diseno_deteccion_picos}),
la digitaliza mediante el ADC integrado del microcontrolador, calcula la
se\~nal de error y la acci\'on de control correspondiente, y comunica el
cambio de temperatura resultante al l\'aser seeder a trav\'es de un
enlace serie \cite{sat2026}. Un conjunto de LEDs y un display indican el
estado del sistema, mientras que la comunicaci\'on con la PC se realiza
por USB \cite{sat2026}.

% PLACEHOLDER: incorporar el diagrama en bloques definitivo de la placa
% (a partir del esquematico de KiCad) y la descripcion de cada bloque
% funcional una vez validada la version final del hardware.

\section{Circuito de comunicaci\'on con el seeder}
\label{sec:comunicacion_seeder}

Para la comunicaci\'on entre el microcontrolador y el l\'aser seeder se
incorpor\'o a la placa un circuito transceptor RS232. Se seleccion\'o el
circuito integrado MAX3232 por su simplicidad de implementaci\'on y su
baja cantidad de componentes externos requeridos
\cite{informeMarzo2026}, conectado a un conector f\'isico DE9
\cite{kicad_hsrl}. El manual de comunicaci\'on del seeder, que detalla
los comandos que deben enviarse por puerto serie para operar el sistema
l\'aser, fue solicitado directamente al fabricante OEM del equipo, NP
Photonics (Estados Unidos) \cite{informeMarzo2026}, y se incluye en el
Anexo~\ref{anexo:manual_seeder}. La correcta funcionalidad de esta
comunicaci\'on fue verificada de acuerdo a lo especificado en dicho
manual \cite{informeAbril2026}.

\section{Dise\~no esquem\'atico y layout en KiCad}
\label{sec:esquematico_layout}

El dise\~no esquem\'atico y el layout de la placa se realizaron en
KiCad~9.0 \cite{informeMarzo2026}, organiz\'andose en hojas
independientes para la etapa de detecci\'on anal\'ogica y para la etapa
digital del sistema \cite{kicad_hsrl}. Entre los componentes principales
incluidos en el dise\~no se encuentran el amplificador AD8004, el
transistor BC817, el transceptor MAX3232, el m\'odulo Raspberry Pi Pico,
el conector coaxial de entrada para la se\~nal del fotodetector y el
conector DE9 para la comunicaci\'on RS232 \cite{kicad_hsrl}.

% PLACEHOLDER: agregar capturas del esquematico y del layout final
% (PCB), asi como la justificacion de la disposicion de capas y el
% stack-up utilizado.

\section{Fabricaci\'on, ensamble y puesta a punto}
\label{sec:fabricacion_ensamble}

La primera revisi\'on de la placa se encarg\'o al fabricante PCBWay,
con un plazo de fabricaci\'on de 48~horas h\'abiles y un env\'io por DHL
de aproximadamente dos semanas \cite{informeMarzo2026}. Una revisi\'on
posterior de la placa se encarg\'o a la empresa JLCPCB
\cite{informeAbril2026}, obteni\'endose el prototipo utilizado para las
validaciones descritas en la Secci\'on~\ref{sec:prototipado_validacion}.

PLACEHOLDER: la fabricaci\'on, el ensamble final y la puesta a punto
integral de la placa de control se encuentran en curso al momento de
redacci\'on de este informe; esta secci\'on debe completarse con la
descripci\'on del proceso de ensamble definitivo, las pruebas de
puesta en marcha y los ajustes realizados sobre la versi\'on final de
la placa.
